% year: תשפ"ג
% semester: ב
% moed: א
% source: exams/מבחנים/תשפג/מועד א תשפג סמסטר ב.pdf
% number: 3ב
% points: 5
% categories: func-analysis

\begin{question}
מצאו את הערכים של פרמטר $b$ עבורם לגרף הפונקציה $\displaystyle f(x)=\frac{1}{x^2+b}$ אין נקודות פיתול
\end{question}

\begin{hint}
חשבו את $f''$ והפרידו למקרים $b>0$, $b=0$, $b<0$. זכרו שנקודת פיתול חייבת להיות בתחום ההגדרה.
\end{hint}

\begin{solution}
נגזור: $f'(x)=-\frac{2x}{(x^2+b)^2}$, ו-
\[ f''(x)=\frac{-2(x^2+b)^2+2x\cdot2(x^2+b)\cdot2x}{(x^2+b)^4}=\frac{6x^2-2b}{(x^2+b)^3}. \]
\begin{itemize}
  \item $b>0$: $f$ מוגדרת על כל $\R$, המכנה חיובי, והמונה $6x^2-2b$ מחליף סימן ב-$x=\pm\sqrt{b/3}$. לכן יש שתי נקודות פיתול.
  \item $b=0$: $f=\frac1{x^2}$, $x\neq0$, ו-$f''=\frac{6}{x^4}>0$: אין נקודות פיתול.
  \item $b<0$: $f$ מוגדרת עבור $x\neq\pm\sqrt{-b}$. המונה $6x^2-2b=6x^2+2|b|>0$ תמיד, ולכן $f''$ לא מתאפסת. סימן $f''$ משתנה רק במעבר דרך $x=\pm\sqrt{-b}$, שאינן בתחום ההגדרה (אסימפטוטות אנכיות), ולכן אינן נקודות פיתול. אין נקודות פיתול.
\end{itemize}
\textbf{תשובה:} לגרף אין נקודות פיתול בדיוק עבור $\boxed{b\le 0}$.
\end{solution}
